\documentclass{article}
\usepackage{/globals/anki}
\ankiprefix{wi}
\ankitopic{WI}
\ankishow
\begin{document}
Die Wirtschaftsinformatik ist ist Kombination aus den Wirtschaftswissenschaften, der Informatik und ein wenig Technik.

Dabei werden an erster Stelle Anwendungssysteme gestaltet, zusammen mit dem Wissen über die Anforderungen und des Umsetzbaren modelliert, und erklärt.

Diese Systeme können in Planungs- und Kontrollsysteme (von NutzerInnen genutzt um entscheidungen zu treffen) und in Administrations- und Dispositionssysteme (\emph{Operative Systeme}, die automatisiert basierend auf Ereignissen handeln) unterteilt werden.

Somit ist es das langzeitige Endziel möglichst große Teile der Wirtschaft zu automatisieren.

\ankinote{ziel}{Ziel}{Vollautomatisierung}
\ankinote{durchsetzung}{Durchsetzung}{Modelle für Anwendungssysteme erstellen}
\ankinote{systeme}{Systeme}{Planungs- und Kontrollsysteme; Operative Systeme (Administrations und Dispositionssysteme)}
\ankinote{pks}{Planungs- und Kontrollsysteme}{Vollautomatisierung}{Entscheidungen mit Daten treffen}
\ankinote{ads}{Administrations und Dispositionssysteme}{Daten generieren}
\end{document}